\section{总结与延伸}

\subsection{核心结论}

\begin{importantbox}{十一条可迁移结论}
\begin{enumerate}
  \item 研究问题会决定状态空间、反馈信号和基础设施形态。
  \item 几何直觉可启动无监督学习，但必须有能修正假设的训练闭环。
  \item Formal proof 的正确性来自 verifier，LLM 负责提出候选和启发搜索。
  \item “Optimal” 必须同时说明 objective、constraint 与 workload。
  \item 小规模正确不能保证大规模数值稳定，训练需要分阶段 canary。
  \item Final run cost 不等于 R\&D cost；失败实验和系统搭建构成冰山主体。
  \item 数据工作不显眼，却常比训练代码更决定模型质量。
  \item Open-weight checkpoint 只是产品栈第一层，不是可用 solution。
  \item 企业定制必须连接数据、评测、部署、监控和回滚。
  \item Reasoning 的可扩展单元是 environment--verifier 闭环，而非更长文本。
  \item Post-training 与开放资产把竞争推向快速迭代、工具、反馈和运营控制。
\end{enumerate}
\end{importantbox}

\subsection{实践作业：设计一个受约束的企业模型系统}

选择一个场景，例如医院文档助手、内部代码 copilot、保险理赔审核或离线 edge agent，完成以下设计：

\begin{enumerate}
  \item 写出不可妥协的约束：数据驻留、延迟、可用性、硬件、责任和更新窗口；
  \item 选择 external API、private cloud、on-prem 或 edge，并说明没有选择其他模式的原因；
  \item 定义 checkpoint-to-solution 的五层组件和 owner；
  \item 设计最小 evaluation suite，区分模型错误、工具错误和系统错误；
  \item 设计 feedback/data loop，说明哪些数据可以回流、保留多久、如何退出；
  \item 给出 base model 升级与 post-training 更新的 canary、rollback 和审计策略。
\end{enumerate}

\begin{knowledgebox}{验收标准}
高质量方案不会只写“使用 Mistral / LLaMA 并微调”。它应能回答：哪个环境产生 ground truth，谁验证结果，数据如何获得，模型如何服务，失败如何定位，何时停止 rollout，以及谁对最终业务结果负责。
\end{knowledgebox}

\subsection{拓展阅读}

\begin{enumerate}
  \item Guillaume Lample et al., \emph{Unsupervised Machine Translation Using Monolingual Corpora Only}, 2018，见 \href{run:source-materials/unsupervised-machine-translation.pdf}{本地 PDF}。
  \item Guillaume Lample et al., \emph{HyperTree Proof Search for Neural Theorem Proving}, 2022，见 \href{run:source-materials/hypertree-proof-search.pdf}{本地 PDF}。
  \item Jordan Hoffmann et al., \emph{Training Compute-Optimal Large Language Models}, 2022，见 \href{run:source-materials/chinchilla-scaling-laws.pdf}{本地 PDF}。
  \item Hugo Touvron et al., \emph{LLaMA: Open and Efficient Foundation Language Models}, 2023，见 \href{run:source-materials/llama.pdf}{本地 PDF}。
  \item Albert Q. Jiang et al., \emph{Mistral 7B}, 2023，见 \href{run:source-materials/mistral-7b.pdf}{本地 PDF}。
  \item Mistral AI 官方 Mistral 7B 发布页：\href{https://mistral.ai/news/announcing-mistral-7b}{Mistral 7B}。
  \item 欧盟委员会 GPAI 指引与时间线：\href{https://digital-strategy.ec.europa.eu/en/policies/guidelines-gpai-providers}{Guidelines for providers of general-purpose AI models}。
\end{enumerate}

\subsection{建议阅读路线}

若目标是复现本讲的推理链，而不是逐篇孤立阅读，可以按以下顺序进行。先读无监督机器翻译论文的 problem setup、训练目标和 ablation，观察一个几何假设如何被 denoising 与 back-translation 修正；再读 HTPS 的环境接口、hyper-tree backup 和 online training，理解 verifier 如何改变搜索算法。随后把 Chinchilla 与 LLaMA 放在一起读，分别写出它们优化的预算约束，并比较小模型长 horizon 训练对推理部署的意义。最后阅读 Mistral 7B 的 GQA/SWA 章节，把 architecture claim 转成 KV cache、带宽、延迟和吞吐的资源账。

\begin{knowledgebox}{阅读记录模板}
每篇材料至少记录五项：\textbf{问题约束}、\textbf{权威反馈信号}、\textbf{核心机制}、\textbf{主要消融或证据}、\textbf{不能推出的结论}。完成后再补一列“如果放进 production，还缺什么系统部件”。这个模板能把论文阅读从结论摘抄转成系统设计练习。
\end{knowledgebox}

\end{document}
